\subsection{Choi Representation II}

\subsubsection{从Choi State得到Kraus Operators}\label{sec:from-choi-to-kraus}

我们之前研究知道通过Vectorization我们可以将一组Kraus Operators映射为一个Choi State。下面我们希望反过来通过Choi State来得到Channel的Kraus Operators。

\bigskip
\hlr{例子：一个特殊channel的choi state}

首先我们计算一个特别的channel的Choi State。我们考虑下面的channel：
\begin{align}
  \mathcal{E}(\rho)=\frac{1}{3}\left(\mathrm{Tr}(\rho)I+\rho^T\right).
\end{align}
我们计算这个channel的Choi State，也就是对于$ \ket{Vec(I)} $进行$ \mathcal{E}(\rho) $的作用我们计算的时候会把前一个和后一个qubit进行区分开来计算：
\begin{align}
  J(\mathcal{E})=\sum_{i,j}\mathcal{E}(|i\rangle\langle j|)\otimes|i\rangle\langle j|
\end{align}
于是我们给出结果是：
\begin{align}
  J(\mathcal{E}) = \frac{1}{3}\left(\sum_iI\otimes|i\rangle\langle i|+\sum_{i, j}\lvert j\rangle\langle i|\otimes|i\rangle\langle j|\right)
\end{align}
我们会发现这个Choi State的第二个部分是一个SWAP Operator。我们下面定义：
\begin{itemize}
  \item SWAP Operator:
    \begin{align}
      SWAP=\sum_{i,j}|i\rangle\langle j|\otimes|j\rangle\langle i| = \sum_{ij} \ketbra{ij}{ji}
    \end{align}
  \item \textbf{对于2 qubit system}，SWAP Operator的本征态，并且本征值为$ (1,1,1,-1) $。我们可以decompose成为下面的形式：
      \begin{align}
        SWAP = | \phi^+ \rangle \langle \phi^+ | + | \phi^- \rangle \langle \phi^- | + | \psi^+ \rangle \langle \psi^+ | - | \psi^- \rangle \langle \psi^- |
      \end{align}
      \textbf{对于更多bit的系统}，我们有：
      \begin{align}
        SWAP = \sum_i \ketbra{ii}{ii} + \sum_{i \neq j} \displaystyle\frac{1}{2} (\ket{ij} + \ket{ji})(\langle ij| + \langle ji|) - \sum_{i \neq j} \displaystyle\frac{1}{2} (\ket{ij} - \ket{ji})(\langle ij| - \langle ji|)
      \end{align}
      其本征值为所有$ \ket{ii} $和$ \displaystyle\frac{1}{\sqrt{2}}(\ket{ij}+\ket{ji}) $对应本征值1，$ \displaystyle\frac{1}{\sqrt{2}}(\ket{ij}-\ket{ji}) $对应本征值-1。
    \item SWAP Operator相当于对于一个2qubit系统进行「转置」也就是：
\begin{align}
  SWAP|ij\rangle=|ji\rangle
\end{align}
\end{itemize}
根据这个信息我们发现Choi State可以写作：
\begin{align}
  J(\mathcal{E})=\frac{1}{3}\left(I\otimes I + SWAP\right) = \frac{2}{3}(|\phi^+\rangle\langle\phi^+|+ |\phi^-\rangle\langle\phi^-| + |\psi^+\rangle\langle\psi^+|)
\end{align}
\rmk{
  这里已经可以看出Choi State是一个半正定矩阵，但是$ \tr $并不是1，所以并不是一个量子态。
}
于是很显然我们发现这个Choi State可以写作三个vectorization的矩阵的形式：
\begin{align}
    A_0=\frac{1}{\sqrt{3}}I,\quad  A_1=\frac{1}{\sqrt{3}}Z,\quad  A_2=\frac{1}{\sqrt{3}}X.
\end{align}
并且我们可以验证这些Kraus Operators满足完备关系：
\begin{align}
  \sum_iA_i^\dagger A_i=\frac{1}{3}{\left(I^2+Z^2+X^2\right)}=\frac{1}{3}(3I)=I.
\end{align}
所以我们成功地通过Choi State得到了Kraus Operators。

\bigskip
\hlr{不唯一的Kraus Operators表示}

但是另一个观察会发现这个Kraus Operators的表示并不是唯一的。观察可以发现等式：
\begin{align}
  \ketbra{00}{00} + \ketbra{11}{11} = \ketbra{\phi^+}{\phi^+} + \ketbra{\phi^-}{\phi^-}
\end{align}
带入上面的Choi State的形式我们发现：
\begin{align}
J(\mathcal{E})=\frac{2}{3}|00\rangle\langle00|+\frac{2}{3}|11\rangle\langle11|+\frac{2}{3}|\psi^+\rangle\langle\psi^+|.  
\end{align}
我们分别验证这个Choi State对应的Kraus Operators：
\begin{align}
  A_0^{\prime}=\sqrt{\frac{2}{3}}|0\rangle\langle0|,\quad A_1^{\prime}=\sqrt{\frac{2}{3}}|1\rangle\langle1|,\quad A_2^{\prime}=\frac{1}{\sqrt{3}}X. 
\end{align}
发现他们也满足完备关系，也就是说他们也是一组合法的Kraus Operators表示。

\bigskip
\hlr{通过Choi State得到Kraus Operators的一般方法}

我们，分析上面的方法发现可以一套完备的方法通过Quantum Channel得到Choi State再得到Kraus Operators：
\begin{enumerate}
  \item \textbf{计算量子态的Choi State}： 通过$ J(\mathcal{E}) = (\mathcal{E} \otimes I)(|Vec(I)\rangle\langle Vec(I)|) $计算Choi State。
  \item \textbf{Choi State的对角化}： 计算Choi State的谱分解$ J(\mathcal{E}) = \sum_i \lambda_i |\psi_i\rangle\langle \psi_i| $。
    \begin{itemize}
      \item 注意：有的时候我们不对角化而是写作一堆纯态的概率混合形式也是可以的。但是我们并不能够保证这样给出的Kraus Operators是完全合法的，需要进行检查！但是对角化的话我们可以保证结果必然是合法的Kraus Operators。
    \end{itemize}
  \item \textbf{得到Kraus Operators}： 通过$ A_i = \sqrt{\lambda_i} \ket{\psi_i}\bra{\psi_i} $得到Kraus Operators。
  \item \textbf{检查完备关系}： 检查$ \sum_i A_i^\dagger A_i = I $是否成立。「检查确实代表了Quantum Channel」
\end{enumerate}

\subsubsection{Kraus Operators的Unitary Mixing Freedom}

\bigskip
\hlr{Unitary Mixing Freedom}

上面已经发现Kraus Operators的表示并不是唯一的。给定一组Kraus Operator我们可以通过一些恒等式给出另外一组Kraus Operator。但有的时候这也不一定能够给出合法的Kraus Operators表示。我们下面给出一个更加一般的结论：
\thm{Unitary Mixing

对于任意的一个半正定的矩阵$ \rho $的两种不同的谱分解：
\begin{align}
  \rho=\sum_k\lambda_k|\lambda_k\rangle\langle\lambda_k|, \quad  \rho=\sum_ip_i|\varphi_i\rangle\langle\varphi_i|.
\end{align}
那么这两组本征态之间存在一个unitary变换$ U $使得：
\begin{align}
  \sqrt{p}_i|\varphi_i\rangle=\sum_kU_{ik}\sqrt{\lambda_k}|\lambda_k\rangle.
\end{align}
}
\rmk{
  这里强调是半正定而非量子态是因为Choi State并不是一个trace为1的量子态。
}
下面我们证明这个结论。首先根据两种谱分解我们给出两种将$ \rho $ purification的方式：
\begin{align}
  |\psi\rangle=\sum_k\sqrt{\lambda_k}|\lambda_k\rangle\otimes|k\rangle,\quad|\varphi\rangle=\sum_i\sqrt{p_i}|\varphi_i\rangle\otimes|i\rangle.
\end{align}
Purification uniqueness定理\cref{thm:purification-uniqueness}告诉我们这两种purification方式之间存在一个unitary变换$ U $使得：
\begin{align}
  |\varphi\rangle=(I\otimes U)|\psi\rangle.
\end{align}
我们写开会发现：
\begin{align}
  \begin{aligned}\sqrt{p_i}|\varphi_i\rangle&\large=(I\otimes\langle i|)(I\otimes U)|\psi\rangle\\&=\sum_k\sqrt{\lambda_k}(I\otimes\langle i|U)(|\lambda_k\rangle\otimes|k\rangle)\\&=\sum_k\sqrt{\lambda_k}|\lambda_k\rangle\cdot(\langle i|U|k\rangle)\\&=\sum_k\sqrt{\lambda_k}U_{ik}|\lambda_k\rangle.\end{aligned}
\end{align}
\qed 


然后我们会发现上面的general量子态的定理如果使用在Choi State上面就得到了我们下面的结论：
\thm{
  Unitary Mixing Freedom for Kraus Operators

  对于任意的一个Quantum Channel $ \mathcal{E} $的两组不同的Kraus Operators表示$ \{A_i\} $和$ \{B_j\} $，那么这两组Kraus Operators之间存在一个unitary变换$ U $使得：
  \begin{align}
    B_j = \sum_i U_{ji} A_i.
  \end{align}
}
其中就是使用了Vectorization线性的性质，可以把外面的unitary拉进去！数学上：
\begin{align}
  |Vec(B_j)\rangle = \sum_i U_{ji} |Vec(A_i)\rangle = \sum_i \ket{Vec(U_{ji} A_i)}.
\end{align}


\subsection{Representation of Quantum Channel}

\hlr{Completely Positive}

我们最后回来讨论一下什么样子的映射才是一个Quantum Channel。这里特别强调一下completely positive的概念。
\begin{itemize}
  \item \textbf{Complete Positive Map}： 一个映射$ \mathcal{E} $是positive map如果对于任意的半正定矩阵$ \rho\geq0 $都有$ (I \otimes \mathcal{E})(\rho)\geq0 $。
\end{itemize}


\bigskip
\hlr{例子：Positive Map but not Completely Positive Map}
  
一个最基本的例子就是转置映射$ T(\rho)=\rho^T $。我们验证这个映射是positive map但是不是completely positive map。
我们考虑转置映射 $ \otimes $ Identity 映射作用在一个bell state上面：
\begin{align}
  (T\otimes I)(|\phi^+\rangle\langle\phi^+|) = \displaystyle\frac{1}{2} SWAP
\end{align}
发现会给出SWAP Operator。我们知道SWAP Operator不是正定的所以这个映射不是completely positive map。
\rmk{
  但是我们会发现$ T \otimes I $作用在某些态上还是正定的映射！只是说存在一些态会被映射成非正定的态。但是这样已经说明了这个map不是completely positive map，也就是说转置是是一个positive map但是不是completely positive map，因而不是一个合法的Quantum Channel。
}

\bigskip 
\hlr{Representation Thm}

\thm{\label{thm:kraus-representation}
  Representation Theorem for Quantum Channel

  一个映射$ \mathcal{E} $满足：
\begin{itemize}
  \item 线性(Linear)
  \item 完全正定(Completely Positive) $  (I \otimes \mathcal{E})(\rho) \geq0 $对于任意$ \rho\geq0 $
  \item 保迹(Trace Preserving) $ \tr(I \otimes \mathcal{E} \rho) = 1 $对于量子态$ \rho $
\end{itemize}
  当且仅当其存在一组Kraus Operators $ \{A_i\} $满足完备关系$ \sum_i A_i^\dagger A_i = I $使得：
  \begin{align}
    \mathcal{E}(\rho) = \sum_i A_i \rho A_i^\dagger 
  \end{align}
}
\textbf{$ \Leftarrow $证明：} 如果一个映射存在Kraus Operators表示，那么我们已经知道这个映射是completely positive map, Linear并且trace preserving。 

\begin{itemize}
  \item 线性是显然的
    \item Trace Preserving:
      \begin{align}
        \tr(\mathcal{E}(\rho)) = \tr\left(\sum_i A_i \rho A_i^\dagger\right) = \tr\left(\rho \sum_i A_i^\dagger A_i\right) = \tr(\rho I) = \tr(\rho).
      \end{align}
    \item Completely Positive: 我们考虑任意量子态$ \ket{\psi} $可以知道：
      \begin{align}\label{eq:completely-positive-proof}
        \bra{\psi}(I \otimes \mathcal{E})(\rho)|\psi\rangle = \sum_i \bra{\psi}(I \otimes A_i) \rho (I \otimes A_i^\dagger)|\psi\rangle  
      \end{align}
      由于$ \rho $是半正定矩阵，我们可以进行对角化：$ \rho = \sum_i \lambda_i \ketbra{i}{i} $。那么上面的式子可以写作：
      \begin{align}
        \bra{\psi}(I \otimes \mathcal{E})(\rho)|\psi\rangle = \sum_i \lambda_i \sum_j |\bra{\psi}(I \otimes A_j)|i\rangle|^2 \geq0.
      \end{align}
      因此我们证明了completely positive。
\end{itemize}

\textbf{$ \Rightarrow $证明：}可以直接使用Choi Representation来证明。根据定理\cref{thm:choi-to-kraus}我们知道必然可以通过Choi State得到Kraus Operators。
当然我们也还有一个更直接的证明方法：

对于任意的Quantum Channel $ \mathcal{E}$，根据Choi Representation的方法必然可以给出一个Choi State $ J(\mathcal{E}) $。根据完全正定的定义我们知道$ J(\mathcal{E}) \geq0 $。因此我们可以对Choi State进行谱分解：
\begin{align}
  J(\mathcal{E}) = \sum_i \lambda_i \ketbra{\lambda_i}{\lambda_i}.
\end{align}
为此我们可以定义一组Operators：
\begin{align}
  A_i \ket{\psi} = \sqrt{\lambda_i} \langle \psi^* |\otimes I | \lambda_i \rangle.
\end{align}
我们注意$ \ket{\lambda_i} $是两个系统的更大的Hilbert Space上的态，而$ \bra{\psi^*} $是作用在辅助系统上面的一个和$ \ket{\psi} $完全一样的态。我们会发现对于任何的一个纯态$ \ket{\psi} $我们有：
\begin{align}
  \sum_k A_k \ketbra{\psi}{\psi} A_k^\dagger = \mathcal{E}(\ketbra{\psi}{\psi}).
\end{align}
由于我们假设了映射是线性的，所以我们可以推广到任意的mixed state $ \rho = \sum_i p_i \ketbra{\psi_i}{\psi_i} $上面：
\begin{align}
  \sum_k A_k \rho A_k^\dagger &= \sum_k A_k \left(\sum_i p_i \ketbra{\psi_i}{\psi_i}\right) A_k^\dagger = \sum_i p_i \left(\sum_k A_k \ketbra{\psi_i}{\psi_i} A_k^\dagger\right) \\&= \sum_i p_i \mathcal{E}(\ketbra{\psi_i}{\psi_i}) = \mathcal{E}\left(\sum_i p_i \ketbra{\psi_i}{\psi_i}\right) = \mathcal{E}(\rho).
\end{align}
进一步的我们考虑Trace Preserving的条件：
\begin{align}
  \tr(\mathcal{E}(\rho)) &= \tr\left(\sum_k A_k \rho A_k^\dagger\right) = \tr\left(\rho \sum_k A_k^\dagger A_k\right) = \tr(\rho).
\end{align}
由于这个等式对于任意的$ \rho $成立，我们知道必然有$ \sum_k A_k^\dagger A_k = I $。因此我们证明了这个映射存在Kraus Operators表示，也就是$ \{A_k\} $是一组合法的Kraus Operators。

\qed 


\bigskip
\hlr{Positivity and Choi Representation}

我们最后讨论一下Choi Representation和Complete Positivity的关系。我们已经知道对于一般的Channel来说Choi State是一个半正定矩阵。我们希望反向也是对的，也就是说如果Choi State是一个半正定矩阵，那么这个映射就是completely positive map。事实上这是正确的：
\thm{
  Choi Representation and Complete Positivity

  一个映射$ \mathcal{E} $是一个Linear,completely positive map当且仅当其满足:
  \begin{align}
    J(\mathcal{E})=\mathcal{E}\otimes I(|vec(I)\rangle\langle vec(I)|) \geq0.
  \end{align}
}
\rmk{注意，我们这里已经不要求$ \mathcal{E} $是quantum channel了，所以不好称呼这个state Choi State。}

\textbf{$ \Rightarrow $证明：} complete positive map的定义已经告诉我们Choi State是一个半正定矩阵。 

\textbf{$ \Leftarrow $证明：} 仔细研究\cref{thm:kraus-representation}的$ \Rightarrow $证明我们会发现。只要一个映射的$ J(\mathcal{E}) $是半正定矩阵我们永远可以找到一组$ \{A_i\} $保证对于任何量子态$ \ket{\psi} $都有：
\begin{align}
  \mathcal{E}(\rho) = \sum_i A_i \rho A_i^\dagger.
\end{align}
那么考虑对于一个半正定矩阵$ \rho\geq0 $我们有：
\begin{align}
  \mathcal{E} \otimes I(\rho) = \sum_i (A_i \otimes I) \rho (A_i^\dagger \otimes I) \geq0.
\end{align}
最后的半正定的证明方法和\cref{eq:completely-positive-proof}的证明是一样的。

